\documentclass[10pt,a4paper]{amsart}
\usepackage[margin=19mm]{geometry}
\usepackage{amsmath,amssymb,mathtools}
\usepackage[scr=rsfs,cal=euler]{mathalfa}
\usepackage{xfrac}
\usepackage[unicode,hidelinks,bookmarks=false]{hyperref}
\newcommand{\ie}{i.e.\ }
\newcommand{\cf}{cf.\ }
\newcommand{\resp}{respectively\ }
\newcommand{\Subsection}{\subsection}
\providecommand{\nobreakdash}{\nobreak\hbox{-}}
\setlength{\emergencystretch}{2em}
\multlinegap0pt
\usepackage{mathtools}
\newtheorem{thm}{Theorem}
\newcommand{\TP}[1]{\textrm{TP}_{#1}}
\DeclareMathOperator{\diag}{diag}
\title{On a connection between total positivity and Bernoulli stopping problems}
\author{Zakaria Derbazi}
\date{Source: arXiv:2411.07103v3 (CC BY 4.0)}
\begin{document}
\maketitle

\noindent\emph{Source and licence.} On a connection between total positivity and Bernoulli stopping problems. Version arXiv:2411.07103v3 (CC BY 4.0), distributed by arXiv under the Creative Commons Attribution 4.0 International licence, http://creativecommons.org/licenses/by/4.0/, as recorded in the arXiv licence metadata for this exact version and shown on the abstract page https://arxiv.org/abs/2411.07103v3. Full licence notice: \texttt{licenses/arxiv-2411.07103-CC-BY-4.0.txt}.\par\smallskip

\noindent\textit{Editorial scope.} Counted unit: the article's thm prop:tp.stochastic.operator (source0.tex lines 241--243). The article's complete original proof of it is delivered with it (source0.tex lines 244--262, one \texttt{proof} environment of 19 lines). No mathematics is written, repaired or re-derived: the statement and the proof are the article's own lines, in the article's own notation, and no retained line is substituted or re-flowed.  No further source-local context is needed: every cross-reference of every retained line resolves inside the two delivered spans.  Nothing was omitted from any retained span, so the package carries no omission record.  No retained line contains a citation, so the wrapper carries no bibliography and no bibliography entry.  No retained line contains a figure environment, an image-inclusion command or drawing code, so no image is delivered and the package declares no asset dependency.  Editorial formatting only, in addition: an AMS \texttt{amsart} preamble that restates the article's own declarations and macros used by the retained lines, namely the environments \texttt{thm} on separate counters, together with the standard packages those lines need.  The retained environments print in this document as Theorem 1 in order of appearance, whereas the article prints the same items with the numbering of its own full document.  The complete native source of the article as distributed in the arXiv e-print for this exact version is bundled unchanged as \texttt{sources/168347/source0.tex}.
	\begin{thm}\label{prop:tp.stochastic.operator}
		The Markov chain $M$  is $\TP{}$.
	\end{thm}

	\begin{proof}
		Multiply each row $i\in S$ of  $P$ by $\prod_{k=1}^{i-1}q_{k}$ and divide  each column   $j \in S$ of the resulting matrix by $ p_j\prod_{k = 1}^{j-1} q_k $, to obtain the matrix $O \coloneqq D_1PD_2^{-1}$, where
		\begin{equation}\label{def.ABC}
			\begin{aligned}
				D_1 &= \diag(1, q_1, q_1q_2. \ldots, \prod_{k=1}^{n}q_{k}),\\
				D_2 &= \diag(1, p_1,  p_2q_1, \ldots,  p_{n}\prod_{k=1}^{n-1}q_{k}),\\[5pt]
				O&\coloneqq D_1PD_2^{-1} = \left(\begin{array}{cccccc}
					0 & 1 & 1 & \cdots& 1\\
					0 & 0 & 1 &   \cdots & 1\\
					\vdots & \ddots  &\ddots &\ddots & \vdots\\
					\vdots & \ddots  &\ddots &\ddots & 1\\		
					0 & 0 & \cdots  &0 &1 \\				
				\end{array}\right),
				%	\end{equation*}
		\end{aligned}
	\end{equation}
	with $n \ge 2$ and $p_\infty = 1$ by convention. To complete the proof, it suffices to demonstrate that \( D_1^{-1} \), \( D_2 \), and \( O \) are \( \TP{} \), since \( P = D_1^{-1} O D_2 \) and the fact that \( \TP{} \) matrices are closed under matrix multiplication. The diagonal matrices \( D_1^{-1} \) and \( D_2 \) are trivially \( \TP{} \). For the matrix $O$, if we exclude the terminal state, the transpose of the remaining submatrix is the kernel representation of the sequence \( (0,1,1,\dots,1)\).  But this  is a constant sequence with a leading zero, hence the corresponding  kernel representation is  \( \TP{} \). Including the terminal state adds a row and column that preserve nonnegativity of all minors. Since all minors of all orders are nonnegative; hence   \( O \) is \( \TP{} \), completing the proof.

\end{proof}

\end{document}
